% concatenated from Frequency.tex
\documentclass[11pt]{amsart}

\usepackage{amsmath,amssymb,latexsym} % ams auxiliaries.
\usepackage{graphicx}
\usepackage{fullpage}
\usepackage[square, numbers]{natbib}   % Arguably a nicer bibliography.
\usepackage{comment}


% \usepackage{times}   % Causes problems with accents & spacing in PDF.





\newtheorem{theorem}{Theorem} % If subsection, get 3.1.2, etc.
\newtheorem{corollary}{Corollary}
\newtheorem{proposition}{Proposition}
\newtheorem{claim}{Claim}
\newtheorem{lemma}{Lemma}
\theoremstyle{definition}
\newtheorem{definition}{Definition}
\newtheorem{example}{Example}
\newtheorem{question}{Question}
\newtheorem{conjecture}{Conjecture}
\newtheorem{remark}{Remark}



\def\Ker{\operatorname{Ker}}
\def\M{\mathsf{M}}
\def\K{\mathsf{K}}
\def\R{\mathbb{R}}
\def\Z{\mathbb{Z}}
\def\T{\mathsf{T}}
\def\F{\mathcal{F}}
\def\aa{\boldsymbol{a}}
\def\bb{\boldsymbol{b}}
\def\cc{\boldsymbol{c}}
\def\ee{\boldsymbol{e}}
\def\qq{\boldsymbol{q}}
\def\rr{\boldsymbol{r}}
\def\ss{\boldsymbol{s}}
\def\uu{\boldsymbol{u}}
\def\vv{\boldsymbol{v}}
\def\xx{\boldsymbol{x}}
\def\yy{\boldsymbol{y}}
\def\zz{\boldsymbol{z}}
\def\llambda{\boldsymbol{\lambda}}
\def\aalpha{\boldsymbol{\alpha}}

\def\zero{\boldsymbol{0}}
\def\supp{\operatorname{supp}}
\def\conv{\operatorname{conv}}
\def\ds{\displaystyle}
\def\KG{\operatorname{KG}}
\def\cd{\operatorname{cd}}
\def\rk{\operatorname{rk}}
\def\alt{\operatorname{alt}}
%\def\H{\mathcal{H}}
\def\id{\operatorname{id}}
\def\diam{\operatorname{diam}}
\def\stab{\operatorname{stab}}
\def\sd{\operatorname{sd}}
\def\mod{\operatorname{mod}}
\def\coind{\operatorname{coind}}

\title{Lower bound for the tau-matching chromatic number via the alternation number}
\author{Saeed Shaebani, Meysam Alishahi}
\date{arxiv:2101.04701v2; CC BY 4.0}
\begin{document}
\maketitle
\noindent\textit{Editorial note.} The counted unit is the article's own Theorem mainthmSaeed at source lines 168-176 with the proof that the author names at source line 192 and opens at source line 194 with the bold header Proof of Theorem mainthmSaeed; the proof environment is not adjacent to the statement because the author's remark at 178-186 and the pointer sentence at 192 lie in between, and the pairing rests on that author-written header, not on a later-first-proof rule.\par\medskip
\noindent\textit{Editorial note.} The retained context at source lines 115-121, 124-126, 137-155, 158-162, 166, 178-186 and 192 is the article's own definition of the tau-matching chromatic number, of the general Kneser hypergraph, of the alternation number, the cited alternation-number bound of Theorem alihaji, the author's announcement of the theorem and the author's pointer to its proof.\par\medskip
\noindent\textit{Editorial note.} The native preamble is retained from source lines 1-3 and 5-69; only the hyperref line at source line 4 is omitted, and no other preamble line is altered.\par\medskip
\noindent\textit{Editorial note.} The bibliography fragment contains verbatim the single bibliography entry cited by the retained lines, namely the Alishahi-Hajiabolhassan paper named in Theorem alihaji.\par\medskip
\noindent\textit{Editorial scope.} One original result: the article's Theorem mainthmSaeed with its complete author proof, plus the article's own definitions of the tau-matching chromatic number, of the general Kneser hypergraph and of the alternation number. The article's Theorem mainthmSaeed2, its sharpness Proposition and the already counted public unit from this source are not counted here and are not used as premises.\par\medskip
\noindent\textit{Allowed standard prerequisites (editorial).}\begin{itemize}
\item Alternation number of a hypergraph. For the cyclic group Z\_r, a vector of nonzero entries is alternating when consecutive nonzero terms differ, alt(X) is the length of the longest alternating subsequence, and alt\_r(H) is the minimum over injective orderings sigma of V(H) of the largest alt(X) for which no colour class spans an edge of H; the article's own definitions are retained at source lines 137-155.. Supplies the invariant |V(H)|-alt\_r(H) that both the counted theorem and its proof use, and the identity alt\_r(H\textasciitilde{}) = alt\_r(H) at source line 218.
\item Cited alternation-number lower bound for the Kneser chromatic number. For every r >= 2 and every hypergraph H the chromatic number of the general Kneser hypergraph of H is at least the ceiling of |V(H)|-alt\_r(H) divided by r-1; the article states this as Theorem alihaji at source lines 159-162 with its own citation to Alishahi and Hajiabolhassan. It is a cited result, not a claim of the counted theorem.. Applied to the augmented hypergraph at source line 214 to get (r-1)|A| >= |V(H\textasciitilde{})|-alt\_r(H\textasciitilde{}).
\item Admissible colour frequencies and matching colourings. An r-colour-frequency mapping assigns to each colour an integer below r; an (A,tau)-matching colouring colours the edges by colours of A so that every nonempty monochromatic matching has size at most the frequency of its colour, and the tau-matching chromatic number is the least size of a finite A admitting such a colouring (source lines 115-121).. Fixes the objects and both conclusions of the counted theorem.
\end{itemize}\medskip
Let ${\mathcal H}$ be a hypergraph with $n$ vertices, and $r$ be a positive integer. Also, let $C$ be an arbitrary nonempty (finite or infinite) set, and regard its elements as our available colors.

\noindent Any arbitrary mapping $\tau: C\longrightarrow \{0,1,\ldots,r-1\}$ is called an {\it $r$-color-frequency mapping}. So, by an $r$-color-frequency mapping, each color in $C$ receives a nonnegative integer less that $r$ as its {\it frequency}. Our aim is coloring the edges of ${\mathcal H}$ by some colors in $C$ in such a way that the size of each nonempty monochromatic matching does not exceed the frequency of the color associated to its edges. Naturally, one is looking for the least possible number of colors needed for such a coloring.

\noindent Formally, for a subset $A$ of $C$, an {\it $(A,\tau)$-matching coloring} of ${\mathcal H}$ is
a coloring of edges of ${\mathcal H}$ by some colors in $A$ in such a way that the size of each nonempty monochromatic matching is bounded above by the frequency of the corresponding color associated to its edges.
The hypergraph ${\mathcal H}$ is called {\it $\tau$-admissible} if ${\mathcal H}$ admits an $(A,\tau)$-matching coloring for some finite subset $A$ of $C$; and in this case, we define the {\it $\tau$-matching chromatic number} of ${\mathcal H}$, denoted by $\chi_M(\tau, {\mathcal H})$, equal to the minimum cardinality of such a finite subset $A$ of $C$. As usual, for the case where ${\mathcal H}$ is not $\tau$-admissible, we set $\chi_M(\tau, {\mathcal H}):=+\infty$.

The {\it general Kneser hypergraph} ${\rm KG}^r({\mathcal H})$ is an $r$-uniform hypergraph which has $E({\mathcal H})$ as its vertex set and  whose edges are formed by $r$ pairwise disjoint edges of ${\mathcal H}$, that is,
$$E\left({\rm KG}^r({\mathcal H})\right) :=\left\{\{e_1,\ldots,e_r\}:\; e_i\cap e_j=\varnothing \mbox{ for all } i\neq j\right\}.$$
It is worth pointing out that the well-known and extensively well-studied chromatic number of general Kneser hypergraph $\chi \left( {\rm KG}^r({\mathcal H}) \right)$ is a special case of $\chi_M(\tau, {\mathcal H})$ where all color frequencies are the fixed constant $r-1$.

For any nonnegative integer $n$, let the symbol $[n]$ denote the set $\{1, 2, \dots , n\}$.
Also, for a positive integer $r\geq 2$,
let $\mathbb{Z}_r=\{\omega,\omega^2,\ldots,\omega^r\}$ be a multiplicative cyclic group of order $r$  with a generator $\omega$ in such a way that
$0 \notin \{\omega,\omega^2,\ldots,\omega^r\}$.

\noindent For a vector $X=(x_1,x_2,\ldots,x_n)\in(\mathbb{Z}_r\cup\{0\})^n\setminus\{(0,0, \dots , 0)\}$, a subsequence of nonzero terms of $X$ is called {\it alternating} if each two consecutive terms of this subsequence are different.  The length of the longest alternating subsequence of $X$ is denoted by ${\rm alt}(X)$. For example,
if we set $r=3$, $n=14$, and $X=(\omega^2 , 0 , 0 , \omega^2,\omega , 0 ,\omega^3 , 0 , \omega^2 , \omega^3 , 0 , \omega^3 , \omega , 0)$, then ${\rm alt}(X)=6$.
Also, we define ${\rm alt}(0,0,\dots,0)$ to be zero.
For an $X=(x_1,x_2,\ldots,x_n)\in(\mathbb{Z}_r\cup\{0\})^n$ and $\epsilon\in \mathbb{Z}_r$, we define $X^\epsilon\subseteq [n]$ to be the set of all indices $j$ such that $x_j=\epsilon$, that is,
$X^\epsilon :=\{j:\; x_j=\epsilon\}$.
By abuse of language, we can write $X=\displaystyle{\left(X^\epsilon\right)}_{\epsilon\in\mathbb{Z}_r}$.

Let ${\mathcal H}$ be a hypergraph with $n$ vertices and $\sigma:[n]\longrightarrow V({\mathcal H})$ be an injective mapping.
We define ${\rm alt}_r({\mathcal H},\sigma)$ to be the maximum possible value of
${\rm alt} (X)$ for  $X=\displaystyle{\left(X^\epsilon\right)}_{\epsilon\in\mathbb{Z}_r}\in (\mathbb{Z}_r\cup\{0\})^n$ such that none of $\sigma(X^\epsilon)$ contains any edge of ${\mathcal H}$.
In other words,
$${\rm alt}_r({\mathcal H},\sigma) :=\max\left\{{\rm alt}(X): X\in(\mathbb{Z}_r\cup\{0\})^n \mbox{ and for each } \epsilon\in \mathbb{Z}_r, \mbox{ we have } E({\mathcal H}[\sigma\left(X^\epsilon\right)])=\varnothing  \right\}.$$
The {\it alternation number of ${\mathcal H}$}, denoted by ${\rm alt}_r({\mathcal H})$, is the minimum possible value for ${\rm alt}_r({\mathcal H},\sigma)$ where the minimum is taken over all
injective mappings $\sigma:[n]\longrightarrow V({\mathcal H})$.

Alishahi and Hajiabolhassan ~\cite{2013arXiv1302.5394A} presented a lower bound for the chromatic number of general Kneser hypergraphs ${\rm KG}^r({\mathcal H})$ in terms of $n$ and ${\rm alt}_r({\mathcal H})$.
\begin{theorem}{\rm (\cite{2013arXiv1302.5394A})}\label{alihaji}
	For an integer $r\geq 2$ and a hypergraph ${\mathcal H}$, we have
	$$\chi\left({\rm KG}^r({\mathcal H})\right)\geq \left\lceil{ |V({\mathcal H})|-{\rm alt}_r({\mathcal H})\over r-1}\right\rceil.$$
\end{theorem}

In the rest of this section, a development of Theorem \ref{alihaji} is established. In this regard, the following theorem provides a lower bound for the $\tau$-matching chromatic number of hypergraphs.

\begin{theorem}\label{mainthmSaeed}
	Let $r\geq 2$ be an integer, $C$ be a nonempty finite or infinite set, and
	$\tau: C \longrightarrow \{0,1,\ldots,r-1\}$ be an $r$-color-frequency  mapping.
	Then, for any $\tau$-admissible hypergraph ${\mathcal H}$ there exists a finite subset
	$B$ of $C$ for which
	$$\sum_{b\in B}\tau(b) \geq  |V({\mathcal H})|-{\rm alt}_r({\mathcal H}) .$$
	Besides, we have
	$$\chi_M(\tau,{\mathcal H})\geq \min \left\{|A|:\;   A \subseteq C, \mbox{ and } A \mbox{ is finite,}  \mbox{ and also } \sum_{a\in A}\tau(a) \geq  |V({\mathcal H})|-{\rm alt}_r({\mathcal H})\right\}.$$
\end{theorem}

\begin{remark}
	The Theorem~\ref{mainthmSaeed} can be viewed as a generalization of Theorem~\ref{alihaji}. To illustrate this, let us consider the set of available colors $C$ to be the set of natural numbers
	$\mathbb{N}$, and define a constant $r$-color-frequency mapping $\tau: \mathbb{N} \longrightarrow \{0,1,\ldots,r-1\}$ by setting $\tau(a) = r-1$ for every $a \in \mathbb{N}$. Under this construction, one obtains:
	$$\chi_M(\tau, {\mathcal H})=\chi\left({\rm KG}^r({\mathcal H})\right)$$
	and also,
	$$\min\left\{|A|:\;   A \subseteq C, \mbox{ and } A \mbox{ is finite,}  \mbox{ and } \sum_{a\in A}\tau(a) \geq  |V({\mathcal H})|-{\rm alt}_r({\mathcal H})\right\}=\left\lceil{ |V({\mathcal H})|-{\rm alt}_r({\mathcal H})\over r-1}\right\rceil.$$
	Therefore, the assertion of Theorem \ref{mainthmSaeed} implies 
	$\chi\left({\rm KG}^r({\mathcal H})\right)\geq \left\lceil{ |V({\mathcal H})|-{\rm alt}_r({\mathcal H})\over r-1}\right\rceil.$
\end{remark}

Having concluded the above remark, we now return to the main argument and present the proof of Theorem \ref{mainthmSaeed}, as follows.

\noindent{\bf Proof of Theorem~\ref{mainthmSaeed}.}
Let us regard an arbitrary finite subset $A$ of $C$ for which ${\mathcal H}$ admits an $(A,\tau)$-matching coloring $c: E({\mathcal H})\longrightarrow A$.
%$\displaystyle\sum_{a\in A}\tau(a) \geq  |V({\mathcal H})|-{\rm alt}_r({\mathcal H})$ follows.
Now, for each $a$ in $A$ with $\tau(a) < r-1$, we add $r-1-\tau(a)$ additional vertices
$x^{(a)}_{1},x^{(a)}_{2},\ldots ,x^{(a)}_{r-1-\tau(a)}$ to ${\mathcal H}$
together with adding all additional singleton edges
$\left\{x^{(a)}_{1}\right\},\left\{x^{(a)}_{2}\right\},\ldots ,\left\{x^{(a)}_{r-1-\tau(a)}\right\}$ to $E({\mathcal H})$
in order to obtain a new hypergraph $\widehat{{\mathcal H}}$.
Also, we extend the edge coloring $c: E({\mathcal H})\longrightarrow A$
to an edge coloring of $\widehat{{\mathcal H}}$, say $\widehat{c}: E\left(\widehat{{\mathcal H}}\right)\longrightarrow A$, in such a way that for each $a$ in $A$ with $\tau(a) < r-1$, all additional singleton edges
$\left\{x^{(a)}_{1}\right\},\left\{x^{(a)}_{2}\right\},\ldots ,\left\{x^{(a)}_{r-1-\tau(a)}\right\}$ are colored by $a$.
On one hand, the function $c: E({\mathcal H})\longrightarrow A$ is an $(A,\tau)$-matching coloring
of ${\mathcal H}$. Hence, for each $a$ in $A$, the size of every matching in ${\mathcal H}$
whose edges are colored by $a$ is less than or equal to $\tau (a)$.
On the other hand,
there are exactly $r-1-\tau(a)$ edges in $E\left(\widehat{{\mathcal H}}\right) \setminus E({\mathcal H})$
that are colored by $a$ under $\widehat{c}: E\left(\widehat{{\mathcal H}}\right)\longrightarrow A$. We conclude that for each $a$ in $A$, the size of every matching in $\widehat{{\mathcal H}}$ whose edges are colored by $a$ is at most $r-1$. Therefore, the function
$\widehat{c}: E\left(\widehat{{\mathcal H}}\right)\longrightarrow A$
is a proper vertex coloring of ${\rm KG}^r\left(\widehat{{\mathcal H}}\right)$. So,
we have
$$|A| \geq \chi\left({\rm KG}^r\left(\widehat{{\mathcal H}}\right)\right)\geq { \left|V\left(\widehat{{\mathcal H}}\right)\right|-{\rm alt}_r\left(\widehat{{\mathcal H}}\right)\over r-1};$$
and therefore, 
$$(r-1)|A| \geq \left|V\left(\widehat{{\mathcal H}}\right)\right|-{\rm alt}_r\left(\widehat{{\mathcal H}}\right).$$
Since 
$(r-1)|A| = \left( \displaystyle\sum_{a\in A}\tau(a) \right) + \left|V\left(\widehat{{\mathcal H}}\right)\right|-|V({\mathcal H})|$ and ${\rm alt}_r\left(\widehat{{\mathcal H}}\right)={\rm alt}_r({\mathcal H})$,
we have
\begin{center}
	$\left( \displaystyle\sum_{a\in A}\tau(a) \right) + \left|V\left(\widehat{{\mathcal H}}\right)\right|-|V({\mathcal H})| \geq \left|V\left(\widehat{{\mathcal H}}\right)\right|-{\rm alt}_r({\mathcal H});$
\end{center}
and therefore,
$$\displaystyle\sum_{a\in A}\tau(a) \geq  |V({\mathcal H})|-{\rm alt}_r({\mathcal H}) .$$
We conclude that for any finite subset $A$ of $C$ for which ${\mathcal{H}}$ admits an $(A,\tau)$-matching coloring, the relation
$\displaystyle\sum_{a\in A}\tau(a) \geq  |V({\mathcal H})|-{\rm alt}_r({\mathcal H}) $ is also satisfied. Accordingly,
$$\chi_M(\tau,{\mathcal H})\geq \min \left\{|A|:\;   A \subseteq C, \mbox{ and } A \mbox{ is finite,}  \mbox{ and also } \sum_{a\in A}\tau(a) \geq  |V({\mathcal H})|-{\rm alt}_r({\mathcal H})\right\} ;$$
as desired.

\hfill$\square$

\begin{thebibliography}{99}
\bibitem{2013arXiv1302.5394A}
M.~Alishahi and H.~Hajiabolhassan.
\newblock On the chromatic number of general {K}neser hypergraphs.
\newblock {\em Journal of Combinatorial Theory, Series B}, 115:186 -- 209,
  2015.
\end{thebibliography}
\end{document}
